%! AP Statistics — Unit 3, Lesson 3.2 — Group Activity
\documentclass[10pt]{article}
\usepackage{apstats-article}
\usepackage{apstats-boxes}

\begin{document}

\pageheader{Unit 3, Lesson 3.2}{Group Activity: Can We Conclude That?}
\namepartnerperiod

\vspace{0.1in}

\begin{scenariobox}[Context]{sky}
\small
For each study below, your group will (1) classify it as an observational study or
an experiment, (2) identify the population and sample, and (3) determine what
conclusions are valid --- and which are not.
\end{scenariobox}

\vspace{0.1in}

\textbf{Study 1:} Researchers asked 1{,}000 randomly selected adults how many hours of
TV they watch per day and recorded their Body Mass Index (BMI). Those who watched
more than 4 hours had significantly higher BMIs.

\begin{practicebox}
\small
\begin{enumerate}[leftmargin=1.5em, itemsep=8pt]
  \item Observational Study \quad / \quad Experiment --- Circle one.
  \item Population: \blank{8cm} \quad Sample: \blank{8cm}
  \item Can we generalize to the population? Why or why not?
        \blank{1cm} \quad \blank{8cm}
  \item Can we conclude that watching TV \emph{causes} higher BMI? Explain.\\
        \writeline\\[1pt]\writeline
  \item Name a possible confounding variable: \blank{8cm}
\end{enumerate}
\end{practicebox}

\vspace{0.1in}

\textbf{Study 2:} Sixty student volunteers at one school were randomly assigned to
either a new reading program or a traditional reading program. After 8 weeks, the
new-program group scored significantly higher on a reading assessment.

\begin{practicebox}
\small
\begin{enumerate}[leftmargin=1.5em, itemsep=8pt]
  \item Observational Study \quad / \quad Experiment --- Circle one.
  \item Was random assignment used? \blank{2cm} \quad Was random selection used? \blank{2cm}
  \item Can we conclude the program \emph{caused} higher scores? Why?\\
        \writeline\\[1pt]\writeline
  \item Can we generalize these results to \emph{all} students? Why or why not?\\
        \writeline\\[1pt]\writeline
\end{enumerate}
\end{practicebox}

\vspace{0.1in}

\textbf{Study 3:} A health department randomly selects 500 adults from its county's
voter registration list and asks each person about their diet and cholesterol levels.
They find a positive association between eating red meat and high cholesterol.

\begin{practicebox}
\small
\begin{enumerate}[leftmargin=1.5em, itemsep=8pt]
  \item Observational Study \quad / \quad Experiment --- Circle one.
  \item Is this a random sample from the county? \blank{3cm}
  \item Can results be generalized to the county? \blank{3cm} To all U.S. adults? \blank{3cm}
  \item Can we conclude red meat \emph{causes} high cholesterol? \blank{3cm}\\
        Justify: \blank{9cm}
\end{enumerate}
\end{practicebox}

\vspace{0.12in}

\begin{practicebox}
\small
\textbf{Group Synthesis:} Fill in the table.

\vspace{4pt}
\renewcommand{\arraystretch}{1.8}
\begin{tabularx}{\linewidth}{@{} >{\bfseries}p{0.22\linewidth} X X @{}}
 & \textbf{Can establish causation?} & \textbf{Can generalize to population?} \\
\hline
Observational study, random sample & & \\
Experiment, random assignment, no random sample & & \\
Experiment, random assignment, random sample & & \\
\end{tabularx}
\end{practicebox}

\end{document}
